\documentclass[english, parskip=full]{scrartcl}
\usepackage[english]{babel}  % Sprache auf Deutsch 
\usepackage[utf8]{inputenc}  % Kodierung der Datei
\usepackage[left=2cm, right=2cm, top=2cm, bottom=3cm, footskip=1.5cm]{geometry}
\usepackage{color}
\usepackage{graphicx, gensymb}
\usepackage{amsfonts, cancel, amssymb}
\usepackage{amsmath, amsthm, array, esint}
\usepackage{bm} % bold math symbols, e.g. \bm{\sigma} etc.
\usepackage{booktabs, multirow}  % schönere Tabellen
%\usepackage{scrlayer-scrpage}    % Manipulation des Seitenstils
%\pagestyle{scrheadings}  % Stil für die Seite setzen
\usepackage{tabularx}
\usepackage{setspace}
%\automark{section}  % Abschnittsnamen als Seitenbeschriftung 
%\setheadsepline{.5pt}    % Kopzeile durch Linie abtrennen
\usepackage[hidelinks]{hyperref}  % Links und weitere PDF-Features
\usepackage{typearea, tikz, pgffor, verbatim}
\usetikzlibrary{calc, arrows.meta,arrows, graphs, quotes, positioning, angles, babel}
\usepackage[cache=false, outputdir=build]{minted}
\usemintedstyle{vs}
\usepackage{placeins} % provides \FloatBarrier

\definecolor{bg}{rgb}{0.9,0.9,0.9}
\newcommand\numberthis{\addtocounter{equation}{1}\tag{\theequation}}
\usepackage{svg}
\usepackage{siunitx}
\usepackage{physics}
\usepackage{tensor}
\usepackage{slashed}
\usepackage{bbm}
\usepackage[compat=1.1.0]{tikz-feynman}

\usepackage{caption}
\captionsetup{
    % width=.8\textwidth,
    % indention=1cm,
    margin=0.5cm,
    format=plain,
    labelfont=bf
}
%\usepackage{subcaption}
%\subcaptionsetup{
%    indention=0cm,
%    format=hang,
%    margin=0.5cm,
%    labelfont=normalfont
%}
\usepackage[english,capitalise]{cleveref}
\usepackage[
    backend=biber,
    sorting=none,
    style=phys,
    giveninits=true,
    sortcites=true,
    citestyle=numeric-comp,
    % url=false,
    % doi=false,
    biblabel=brackets,
    % eprint=false,
    maxcitenames=3]{biblatex}
\usepackage{csquotes}
%\addbibresource{bibliography.bib}

\usepackage{mathtools}
% use \abs for absolute values
% use \abs for \left ... \right version and \abs[\bigg for \bigg version etc.
\let\abs\undefined
\let\bra\undefined
\let\braket\undefined
\DeclarePairedDelimiter\abs{\vert}{\vert}
\DeclarePairedDelimiter\kla{(}{)}
\DeclarePairedDelimiter\klb{[}{]}
\DeclarePairedDelimiter\klc{\{}{\}}
\DeclarePairedDelimiter\bra{\langle}{\rangle}
\DeclarePairedDelimiterX\brb[2]{\langle #1}{\rangle}{\delimsize\vert #2 }
\DeclarePairedDelimiterX\brc[3]{\langle #1}{\rangle}{\delimsize\vert #2 \delimsize\vert #3 }

\renewcommand{\i}{\text{i}}
\newcommand{\e}{\text{e}}
\DeclareMathOperator{\sgn}{sgn}
\newcommand{\kom}[1]{\overset{\substack{#1 \\ |}}{=}}
\newcommand{\koml}[2]{\overset{\substack{#1 \\ |}}{#2}}
\newcommand{\intx}[4]{\int_{#1}^{#2} #3 \,\mathrm{d}#4}
\newcommand{\intr}[2]{\int_{\mathbb{R}} #1 \, \mathrm{d}#2}
\newcommand{\intrnd}[3]{\int_{\mathbb{R}^{#1}} #2 \, \mathrm{d}^{#1}#3}
\newcommand{\intrpi}[2]{\int_{\mathbb{R}} #1 \, \frac{\mathrm{d}#2}{2\pi}}
\newcommand{\intrndpi}[3]{\int_{\mathbb{R}^{#1}} #2 \, \frac{\mathrm{d}^{#1}#3}{(2\pi)^{#1}}}
\newcommand{\oper}[2]{\vert #1 \rangle \langle #2 \vert}
\newcommand{\Text}[1]{\text{#1}}    % this is relevant because \text does not autocomplete to the default '\text{@}' but rather '\textbf' etc.
\newcommand{\powe}[1]{\e^{#1}}
\newcommand{\pow}[2]{{#1}^{#2}}
\newcommand{\Ket}[1]{\vert #1 \rangle}
\newcommand{\Bra}[1]{\langle #1 \vert}
\newcommand*{\Fourier}{\textsc{Fourier}\ }
\newcommand*{\F}{\mathcal{F}}
\DeclareMathOperator{\arsinh}{arsinh}
\DeclareMathOperator{\arcosh}{arcosh}
\DeclareMathOperator{\artanh}{artanh}
\DeclareMathOperator{\sinc}{sinc}
\DeclareMathOperator{\cscc}{cscc}
\DeclareMathOperator{\sinhc}{sinhc}
\DeclareMathOperator{\cschc}{cschc}
\newcommand{\conv}{\mathop{\scalebox{3.5}{\raisebox{-0.38ex}{\makebox[0.5\width][c]{$\ast$}}}}\limits}
\newcommand*{\unity}{\text{\usefont{U}{bbold}{m}{n}1}}   % operator one from :: https://tex.stackexchange.com/questions/566780/import-mathbb1-character-from-the-unicode-math-package

\renewcommand{\d}{\text{d}}
\renewcommand{\r}{\text{r}}
\renewcommand{\t}{\text{t}}
\renewcommand{\a}{\text{a}}
\renewcommand{\o}{\text{o}}
\renewcommand{\F}{\text{F}}
\renewcommand{\c}{\text{c}}
\newcommand{\A}{\text{A}}
\newcommand{\B}{\text{B}}
\newcommand{\R}{\text{R}}
\newcommand{\M}{\text{M}}
\newcommand{\m}{\text{m}}
\newcommand{\s}{\text{s}}
\newcommand{\f}{\text{f}}
\newcommand{\K}{\text{K}}
\newcommand{\hc}{\text{h.c.}}
\newcommand{\kf}{k_\F}
\newcommand{\vf}{v_\F}
\newcommand{\const}{\text{const.}}
\renewcommand{\L}{\text{L}}
\newcommand{\gray}[1]{\bgroup\color{white!40!black}#1\egroup}
\newcommand{\TODO}[1]{\bgroup\color{red!60!black}#1\egroup}
\newcommand\QnA[1]{\bgroup\color{blue}#1\egroup}
\newcommand\highlight[1]{\bgroup\color{green!60!black}#1\egroup}

\newcommand{\cabx}[3]{\begin{smallmatrix}\gray{A}&\gray{:}&#1 \\ \gray{B}&\gray{:}&#2 \\ \gray{\Xi}&\gray{:}&#3\end{smallmatrix}}
\newcommand{\zabx}[3]{\begin{smallmatrix}\gray{A}&\gray{:}&#1 \\ \gray{B}&\gray{:}&#2 \\ \gray{Z}&\gray{:}&#3\end{smallmatrix}}
\newcommand{\abx}[3]{\begin{smallmatrix}#1 \\ #2 \\ #3\end{smallmatrix}}

\newcommand{\normord}[1]{
  {:\mathrel{\mspace{1mu}#1\mspace{1mu}}:}
}
\newcommand{\varnormord}[1]{
  {\mathrel{\mspace{1mu}\substack{\ast\\[-1.5pt] \ast}#1\substack{\ast\\[-1.5pt] \ast}\mspace{1mu}}}
}

% punctuation styles (see https://tex.stackexchange.com/questions/56063/how-to-properly-typeset-footnotes-superscripts-after-punctuation-marks)
\let\origfootnote\footnote
\renewcommand{\footnote}[1]{\kern.06em\origfootnote{#1}}
\newcommand{\punctfootnote}[1]{\kern-.06em\origfootnote{#1}}

% Multiline equations
\newcommand{\bsplit}[1]{\begin{aligned}[t]#1\end{aligned}}

\newcommand*{\titel}{Path Integral Calculations}
\newcommand*{\autor}{Joachim Schwardt}

\hypersetup{pdfauthor={\autor}, pdftitle={\titel}}

%\titlehead{titlehead}
\subject{Quantum Physics in 1D}
\title{\titel}
\author{\autor}
\date{\today}

\begin{document}

%\begin{titlepage}
\maketitle  % Titel setzen
%\tableofcontents  % Inhaltsverzeichnis setzen
%\end{titlepage}


\section{Quadratic Luttinger Hamiltonian}
We want to compute the basic correlation functions for a quadratic Luttinger liquid described by
\begin{align}
H &= \frac{u}{2\pi} \intr{\klb*{K\kla*{\partial_x\theta}^2 + \frac{1}{K} \kla*{\partial_x\phi}^2}}{x}.
\end{align}
In order to use the path integral formalism to evaluate expectation values we note that $\Pi = \frac{1}{\pi}\partial_x\theta$ is the conjugate momentum of $\phi$ and switch to a Lagrangian description with the action
\begin{align}
S &= \intx{}{}{L(\phi,\partial_t\phi, \partial_x\phi)}{t} = \intx{}{}{\klb*{\Pi\partial_t \phi - H}}{t}.
\end{align}
%Performing the so called Wick rotation $t=-\i\tau$ to imaginary time $\tau$ and writing $r=(x,\tau u)$ we find
%\begin{align}
%S &= \intx{0}{\beta}{\intr{\klb*{\i\frac{1}{\pi}\nabla\theta \partial_\tau\phi - \frac{uK}{2\pi} (\nabla\theta)^2 - \frac{u}{2\pi K}(\nabla\phi)^2 + J(r)\phi(r)}}{x} }{\tau}
%\end{align}
%where we added a source term $J\nabla\phi$ to later evaluate expectation values involving combinations of $\phi$ and $\nabla\phi$.
%At this point we perform a Fourier transformation
%\begin{align}
%F(r) &= \frac{1}{2\pi\beta} \sum_{\omega_n^\B} \intr{\powe{-\i\omega_n^\B \tau + \i kx} F\kla*{\omega_n^\B,k}}{k}
%\end{align}
%for the fields $F\in\klc*{\phi,\theta}$.
Since $\partial_t\phi = \frac{\delta H}{\delta \Pi} =\pi u K \Pi$ we can rewrite the action as
\begin{align}
S &= \intx{ }{ }{\intr{\klb*{\frac{1}{\pi uK} (\partial_t\phi)^2 - \frac{uK}{2\pi} \kla*{\frac{\partial_t\phi}{uK}}^2 - \frac{u}{2\pi K} (\partial_x\phi)^2}}{x} }{t}\\
&= \frac{1}{2\pi uK} \intx{ }{ }{\intr{\klb*{(\partial_t\phi)^2 - u^2 (\partial_x\phi)^2}}{x} }{t}.
\end{align}
Performing the Wick rotation $t=-\i\tau$ to imaginary time $\tau$ and writing $r=(x,u\tau)$ gives
\begin{align}
S &= -\frac{u}{2\pi K} \intx{0}{\beta}{\intr{\klb*{\frac{1}{u^2} (\partial_\tau\phi)^2 + (\partial_x\phi)^2 }}{x} }{\tau}.
\end{align}
With $y=u\tau$ and partial integration this becomes
\begin{align}
S &= \frac{1}{2\pi K} \int \phi \nabla^2 \phi\,\d^2r.
\end{align}
The partition function can be evaluated as
\begin{align}
Z[J] &= \int \powe{-S[\phi] + \int J(r)\phi\,\d^2r}\,\mathcal{D}\phi = \sqrt{\frac{2\pi K}{\det (\nabla^2)}} \powe{\frac{\pi K}{2} \int J(r) D(r,r') J(r')\,\d^2r \,\d^2r'}
\end{align}
%%% https://physics.stackexchange.com/questions/605653/determinant-of-differential-operator-partial2-m2
with the Kernel satisfying the differential equation
\begin{align}
\nabla^2 D(r,r') &= \delta(r-r').
\end{align}
The evaluation of the functional determinant is not easy due to the need for regularization and the treatment of the zero-modes.
However, since the term is independent of the source $J$ it will drop out in any calculation of correlators.
For example
\begin{align}
\bra*{\phi(r)\phi(0)} &= \frac{1}{Z[J]} \frac{\delta^2}{\delta J(r) \delta J(0)} Z[J] \Big\vert_{J=0} = \pi K D(r,0).
\end{align}
Let us take a moment to explicitly evaluate $D$.
Translational invariance implies $D(r,r')= D(r-r')$ and thus we only need to solve the equation $\nabla^2D(r) = \delta(r)$.
Note that this reduces to the typical Green's function of the Laplace equation if we do not impose boundary conditions.
In our case however, we do have periodic boundary conditions due to the requirement $\phi(\tau+\beta) = \phi(\tau)$.
Writing $b=\beta u$ this translates to $D(y+b)=D(y)$.
To build this condition into our ansatz we use the Fourier expansion
\begin{align}
D(x,y) &= \frac{1}{2\pi b} \sum_{n\in\mathbb{Z}} \intr{\powe{\i kx - \i w_n y} D(k,w_n)}{k}
\end{align}
with the discrete frequencies $w_n = \frac{2\pi n}{b}$.
Using the identity
\begin{align}
\frac{1}{2\pi b} \intx{0}{b}{\intr{\powe{\i kx - \i w_n y}}{x} }{y} &= \delta(k) \delta_{n,0}
\end{align}
allows us to write
\begin{align}
D(k,w_n) &= \intx{0}{b}{\intr{\powe{\i w_n y - \i kx} D(x,y)}{x} }{y}
\end{align}
which we may use to transform the differential equation to
\begin{align}
1 &= -\kla*{w_n^2 + k^2} D(k,w_n).
\end{align}
Thus we find
\begin{align}
D(x,y) &= \frac{1}{2\pi b} \sum_{n\in\mathbb{Z}} \powe{-\i w_n y} \intr{\powe{\i kx} \frac{-1}{w_n^2 + k^2} \powe{-\alpha |k|}}{k} \\
&= D(0,0) + \frac{1}{2\pi b} \sum_{n\in\mathbb{Z}} \intr{\kla*{1 - \powe{\i kx - \i w_ny}} \frac{-1}{w_n^2 + k^2} \powe{-\alpha |k|}}{k}\\
&= D(0,0) + \frac{1}{2b} \sum_{n\in\mathbb{Z}} \kla*{1 - \powe{- |w_n x| - \i w_ny}} \powe{-\alpha |w_n|} \frac{1}{|w_n|}
\end{align}
Here we introduced a short-distance cut-off $\alpha$ to ensure convergence and split off the divergent constant.
This procedure has the simple interpretation that the correlator we really care about is
\begin{align}
\bra*{\klb*{\phi(r) - \phi(0)}^2} &= \bra*{\phi(r)^2} + \bra*{\phi(0)^2} - 2\bra*{\phi(r) \phi(0)} = -2\bra*{\phi(r)\phi(0) - \phi(0)^2} \\
&= -2\pi K \klb*{D(r) - D(0)}.
\end{align}
Thus we rewrite in terms of this difference and find\footnote{$\sum_{n>0} \frac{1}{n} c^n = \int \frac{1}{1-c}\,\d c = -\log(1-c)$.}
\begin{align}
D(x,y) - D(0,0) &= -\frac{|x|}{2b} + \frac{1}{b} \Re \sum_{n>0} \frac{1}{w_n} \kla*{1 - \powe{- w_n |x| - \i w_ny}} \powe{-\alpha w_n}\\
&= -\frac{|x|}{2b} + \frac{1}{2\pi} \log\kla*{1 - \powe{-\alpha \frac{2\pi}{b}}} - \frac{1}{2\pi} \Re \log\kla*{1 - \powe{-\frac{2\pi}{b} (|x| + \i y + \alpha)}}\\
&= -\frac{1}{2\pi} \Re \log \kla*{\powe{\frac{\pi\alpha}{b} + \frac{\pi |x|}{b}} \frac{1 - \powe{-\frac{2\pi}{b}(x+\i y+\alpha)}}{2\sinh\kla*{\frac{\pi\alpha}{b}}} \powe{\i\frac{\pi}{b}y} } \\
&\simeq \frac{1}{2\pi} \Re \log f_-(r) + \frac{1}{2\pi} \log \kla*{\frac{\pi\alpha}{b}}
\end{align}
where we introduced $f_\ell(r) = \csc\kla*{\frac{\pi}{b} \klb*{y + \i\ell x}}$ and assumed $\alpha \ll r$.
%%% integrate sin(by) cos(by) b /(sin^2(by) + sinh^2(bx)) dx
%xv=0.76;yv=np.pi
%subs={x:xv,y:yv}
%exact=complex(-1j*sy.arg(np.tan(yv) + 1j * np.tanh(xv)).evalf())
%#numeric=-1j *np.sum([integrate.quad(lambda k: 2*n/k * np.sin(k*xv)*np.sin(2*n*yv) / (k**2 + 4*n**2), -np.inf, np.inf, limit=100)[0] for n in range(-100, 101) if n != 0])
%#analytic1 = 1j*np.sign(xv) * np.log((np.exp(-1j*yv) - np.exp(1j*yv)) / (np.exp(-1j*yv) - np.exp(1j*yv) * np.exp(-2*np.abs(xv)))).imag
%res = sy.integrate(sy.diff(-sy.log(sy.sin(y)**2 + sy.sinh(x)**2) / 4, y) * sy.I, (x, 0, xv))
%zr = np.exp(complex(res.subs(subs).evalf()))
%ze = np.exp(complex(-1j*sy.arg(np.tan(yv) + 1j * np.tanh(xv)).evalf()))
%plt.plot([zr.real, ze.real], [zr.imag, ze.imag], ls='', marker='o')
%plt.plot(np.cos(t:=np.linspace(0, 2*np.pi, 300)), np.sin(t))
To evaluate correlators containing $\theta$ we use the relation $\partial_x\theta = \frac{1}{uK} \partial_t\phi = \frac{\i}{K} \partial_y\phi$.
For example, this leads to
\begin{align}
\bra*{\klb*{\theta(r) - \theta(0)}^2} &= -2\bra*{\theta(r)\theta(0) - \theta(0)^2} = -2\intx{0}{x}{\bra*{\partial_{x'}\theta(x',y) \theta(0,0)}}{x'}\\
&= -\frac{2\i}{K} \partial_y \intx{0}{x}{\intx{-\infty}{0}{\bra*{\phi(x',y) \partial_{x''}\theta(x'',y'')} \Big\vert_{y''=0}}{x''}}{x'}\\
&= \frac{2}{K^2} \partial_y\partial_{y''} \intx{0}{x}{\intx{-\infty}{0}{\bra*{\phi(x'-x'',y-y'')\phi(0,0)} \Big\vert_{y''=0}}{x''}}{x'}\\
&= -\frac{2}{K^2} \partial_y^2 \intx{0}{x}{\intx{-\infty}{0}{\bra*{\phi(x'-x'', y)\phi(0,0)}}{x''}}{x'}\\
&= \frac{2}{K^2} \intx{0}{x}{\intx{-\infty}{0}{\bra*{\partial_{x'}^2 \phi(x'-x'', y)\phi(0,0)}}{x''}}{x'}\\
&= -\frac{2}{K^2} \intx{0}{x}{\bra*{\partial_{x'} \phi(x', y)\phi(0,0)}}{x'}\\
&= -\frac{2}{K^2} \bra*{\phi(r)\phi(0) - \phi(0)^2} = \frac{1}{K^2} \bra*{\klb*{\phi(r) - \phi(0)}^2} = \frac{1}{K}F_1(r).
\end{align}
We can also use similar manipulations to show that\footnote{$\arg(z) + \alpha = \arg(z) + \arg \powe{\i \alpha} = \arg \klb*{z \powe{\i \alpha}}$}
\begin{align}
F_2(r) &\equiv 2\bra*{\theta(r) \phi(0) - \theta(0)\phi(0)} = 2\intx{0}{x}{\bra*{\partial_{x'} \theta(x',y)\phi(0,0)}}{x'}\\
&= \frac{2\i}{K} \intx{0}{x}{\partial_y \bra*{\phi(x',y)\phi(0,0)}}{x'} = -\frac{\i}{K} \intx{0}{x}{\partial_y \bra*{\klb*{\phi(x',y) - \phi(0,0)}^2}}{x'}\\
&= -\i \intx{0}{x}{\partial_y F_1(x',y)}{x'} = -\i \intx{0}{x}{ \frac{1}{2}\partial_y \log\klb*{\sin^2\kla*{\frac{\pi y}{b}} + \sinh^2\kla*{\frac{\pi x'}{b}}} }{x'}\\
&= -\i \intx{0}{x}{\frac{\sin\kla*{\frac{\pi y}{b}} \cos\kla*{\frac{\pi y}{b}} \frac{\pi}{b}}{\sin^2\kla*{\frac{\pi y}{b}} + \sinh^2\kla*{\frac{\pi x'}{b}}}}{x'} = \i \arctan\klb*{\tan\kla*{\frac{\pi y}{b}} \coth\kla*{\frac{\pi x'}{b}}}\bigg\vert_{x'=0}^{x'=x}\\
&= \i \arg\klb*{\tanh\kla*{\frac{\pi x}{b}} + \i \tan\kla*{\frac{\pi y}{b}}} - \frac{\pi \i}{2}\\
&= \i \arg\klb*{-\i \tanh\kla*{\frac{\pi x}{b}} + \tan\kla*{\frac{\pi y}{b}}} = -\i \arg\klb*{\tan\kla*{\frac{\pi y}{b}} + \i \tanh\kla*{\frac{\pi x}{b}}}.
\end{align}
BEGIN NEW\\
Note that using $\klb*{\phi(x), \theta(y)} = \i \frac{\pi}{2} \sgn(x-y)$ we find
\begin{align}
\partial_t \phi &= \frac{1}{\i} \klb*{\phi, H} = \frac{u}{2\pi\i} \intr{K \klb*{\phi(x), \kla*{\partial_y\theta(y)}^2}}{y} = \frac{uK}{\pi\i} \intr{\partial_y\theta(y) \partial_y \klb*{\phi(x), \theta(y)}}{y} = uK\partial_x\theta,\\
\partial_t\theta &= \frac{1}{\i} \klb*{\theta, H} = \frac{u}{2\pi\i} \intr{\frac{1}{K} \klb*{\theta(x), \kla*{\partial_y\phi(y)}^2}}{y} = \frac{u}{\pi\i K} \intr{\partial_y\phi \partial_y\klb*{\theta(x), \phi(y)}}{y} = -\frac{u}{K} \partial_x\phi.
\end{align}
Replacing $t=-\i \tau = \frac{y}{\i u}$ yields $\partial_x\theta = \frac{\i}{K} \partial_y\phi$ and $\partial_y\theta = \frac{\i}{K} \partial_x\phi$.
This allows us to represent $\theta$ entirely in terms of $\phi$ up to a constant via
\begin{align}
\theta(x,y) &= \theta(0,y) + \intx{0}{x}{\partial_{x'} \theta(x',y)}{x'} = \theta(0,0) + \frac{\i}{K} \intx{0}{y}{\partial_{x} \phi(0,y')}{y'} + \frac{\i}{K} \intx{0}{x}{\partial_y\phi(x',y)}{x'}.
\end{align}
This is useful for computing expectation values involving $\theta$ and for example yields
\begin{align}
2\bra*{\theta(r)\phi(0) - \theta(0)\phi(0)} &= \frac{2\i}{K} \intx{0}{y}{\partial_x \bra*{\phi(x,y')\phi(0,0)}}{y'}\Big\vert_{x=0} + \frac{2\i}{K} \intx{0}{x}{\partial_y\bra*{\phi(x',y)\phi(0,0)}}{x'}\\
&= 2\i\pi \intx{0}{y}{\partial_x D(x,y') \Big\vert_{x=0}}{y'} + 2\i\pi \intx{0}{x}{\partial_y D(x',y)}{x'}\\
&= \pi\i \intx{0}{x}{\frac{-1}{2\pi} \frac{2\sin\kla*{\frac{\pi y}{b}} \cos\kla*{\frac{\pi y}{b}} \frac{\pi}{b}}{\sin^2\kla*{\frac{\pi y}{b}} + \sinh^2\kla*{\frac{\pi x'}{b}}}}{x'}\\
&= \i \arctan\klb*{\tan\kla*{\frac{\pi y}{b}} \coth\kla*{\frac{\pi x'}{b}}}\bigg\vert_{x'=0}^{x'=x}\\
&= \i \arg\klb*{\tanh\kla*{\frac{\pi x}{b}} + \i \tan\kla*{\frac{\pi y}{b}}} - \frac{\pi \i}{2}\\
&= \i \arg\klb*{-\i \tanh\kla*{\frac{\pi x}{b}} + \tan\kla*{\frac{\pi y}{b}}} = -\i \arg\klb*{\tan\kla*{\frac{\pi y}{b}} + \i \tanh\kla*{\frac{\pi x}{b}}}.
\end{align}
Note that the related expectation value $\bra*{\phi(0)\theta(r) - \phi(0)\theta(0)}$ corresponds to
\begin{align}
\frac{\i}{K} \intx{0}{x}{\partial_y \bra*{\phi(0,0)\phi(x',y)}}{x'} &= \i\pi \intx{0}{x}{\partial_y D(-x',-y)}{x'} = -\i\pi \intx{0}{x}{\partial_y D(x',y)}{x'} = -?!??
\end{align}
Similarly we can evaluate
\begin{align}
\bra*{\klb*{\theta(r) - \theta(0)}^2} &= -2\bra*{\theta(r) \theta(0) - \theta(0)^2}\\
&= -\frac{2\i}{K} \intx{0}{y}{\partial_x \bra*{\phi(x,y')\theta(0,0)}\Big\vert_{x=0}}{y'} - \frac{2\i}{K} \intx{0}{x}{\partial_y \bra*{\phi(x',y) \theta(0,0)}}{x'}\\
&= -\frac{2\i}{K} \intx{0}{y}{\partial_x \bra*{\phi(x,0)\theta(0,-y')}\Big\vert_{x=0}}{y'} - \frac{2\i}{K} \intx{0}{x}{\partial_y \bra*{\phi(0,y) \theta(-x',0)}}{x'}\\
&= \bsplit{&-\frac{2\i}{K} \intx{0}{y}{\partial_x \intx{0}{-y'}{\partial_{y''} \bra*{\phi(x,0)\theta(0,y'')}\Big\vert_{x=0} }{y''} }{y'} \\
&- \frac{2\i}{K} \intx{0}{x}{\partial_y \intx{0}{-x'}{\partial_{x''} \bra*{\phi(0,y)\theta(x'',0)}}{x''}}{x'}}\\
&= \bsplit{&\frac{2}{K^2} \intx{0}{y}{\partial_x \intx{0}{-y'}{\partial_{x'} \bra*{\phi(x,0)\phi(x',y'')}\Big\vert_{x,x'=0} }{y''} }{y'} \\
&+ \frac{2}{K^2} \intx{0}{x}{\partial_y \intx{0}{-x'}{\partial_{y'} \bra*{\phi(0,y)\phi(x'',y')} \Big\vert_{y'=0}}{x''}}{x'}}\\
&= \bsplit{&-\frac{2}{K^2} \intx{0}{y}{\intx{0}{-y'}{\partial_x^2 \bra*{\phi(x-x',-y'')\phi(0,0)} \Big\vert_{x,x'=0}}{y''}}{y'}\\
&- \frac{2}{K^2} \intx{0}{x}{\intx{0}{-x'}{\partial_y^2 \bra*{\phi(-x'',y-y')\phi(0,0)} \Big\vert_{y'=0}}{x''}}{x'}}\\
&= \bsplit{&-\frac{2}{K^2} \intx{0}{y}{\intx{0}{y'}{\partial_{y''}^2 \bra*{\phi(0,y'')\phi(0,0)}}{y''}}{y'}\\
&- \frac{2}{K^2} \intx{0}{x}{\intx{0}{x'}{\partial_{x''}^2 \bra*{\phi(x'',y)\phi(0,0)}}{x''}}{x'}}\\
&= -\frac{2}{K^2} \bra*{\phi(0,y)\phi(0,0) - \phi(0,0)^2} - \frac{2}{K^2} \bra*{\phi(x,y)\phi(0,0) - \phi(0,0)^2}\\
&= 
\end{align}
\\END NEW\\
A slightly different approach can be taken to evaluate the expectation values
\begin{align}
\bra*{\prod_j \powe{\i A_j \phi(r_j) + \i B_j\theta(r_j)}} &= \frac{1}{Z_0} \int \powe{-S[\phi]} \prod_j \powe{\i A_j \phi(r_j) + \i B_j \frac{\i}{K} \partial_{y_j} \intx{-\infty}{x_j}{\phi(x_j', y_j)}{x_j'}}\,\mathcal{D}\phi\\
&= \powe{\frac{\pi K}{2}\int J(r) D(r-r') J(r')\,\d^2r\,\d^2r' }
\end{align}
where we introduced $J(r) = \sum_j \i A_j \delta(r-r_j) + \i\sum_j B_j \frac{-\i}{K}  $ WIP!!!
\subsection{General expectation value identity}
Let $C_j$ be bosonic operators with $\mathbb{C}$-valued commutators.
Then repeated application of Baker-Campbell-Hausdorff yields
%%% this follows BOSONIZATION FOR BEGINNERS (94) and others
\begin{align}
\prod_j \powe{C_j} &= \powe{C_1} \powe{C_2} \dots \powe{C_n} = \powe{C_1+C_2}\dots \powe{C_n} \powe{\frac{1}{2} \klb*{C_1,C_2}} = \dots = \powe{\sum_j C_j} \powe{\frac{1}{2} \sum_{j<k} \klb*{C_j,C_k}}.
\end{align}
Since we assumed the commutators to be simple numbers we can make the trivial replacement $\klb*{C_j,C_k} = \bra*{\klb*{C_j,C_k}}$.
Since the operators are bosonic we find the expectation value
\begin{align}
\bra*{\prod_j \powe{C_j}} &= \powe{\frac{1}{2} \bra*{\kla*{\sum_j C_j}^2}} \powe{\frac{1}{2} \sum_{j<k} \klb*{C_j,C_k}}\\
&= \powe{\frac{1}{2} \sum_j \bra*{C_j^2} + \frac{1}{2} \sum_{j<k} \bra*{C_jC_k} + \frac{1}{2} \sum_{j<k} \bra*{C_kC_j}} \powe{\frac{1}{2} \sum_{j<k} \bra*{\klb*{C_j,C_k}}}\\
&= \powe{\frac{1}{2} \sum_j \bra*{C_j^2}} \powe{\sum_{j<k} \bra*{C_jC_k}} \equiv \powe{S_1} \powe{S_2}.
\end{align}
For the special case of $C_j = \i A_j\phi(r_j) + \i B_j \theta(r_j)$ this yields
\begin{align}
\powe{S_1} &= \powe{-\frac{1}{2} \sum_j \klb*{A_j^2 \bra*{\phi(0)^2} + B_j^2 \bra*{\theta(0)}^2 + A_jB_j \bra*{\phi(0)\theta(0) + \theta(0)\phi(0)}}},\\
\powe{S_2} &= \powe{-\sum_{j<k} \klb*{A_jA_k \bra*{\phi(r_j-r_k)\phi(0)} + B_jB_k \bra*{\theta(r_j-r_k)\theta(0)} + A_jB_k \bra*{\phi(r_j-r_k)\theta(0)} + B_jA_k \bra*{\theta(r_j-r_k)\phi(0)}}}
\end{align}
Writing $F_1(r) = \bra*{\klb*{\phi(r) - \phi(0)}^2}$, $F_2(r) = 2\bra*{\theta(r)\phi(0) - \theta(0)\phi(0)}$, $F_3(r) = \bra*{\klb*{\theta(r) - \theta(0)}^2}$ and $F_4(r) = 2\bra*{\phi(r)\theta(0) - \phi(0)\theta(0)}$ yields
\begin{align}
\powe{S_2} &= \powe{\frac{1}{2} \sum_{j<k} \klb*{A_jA_k F_1(r_j-r_k) + B_jB_kF_3(r_j-r_k) - A_jB_k F_2(r_j-r_k) - A_kB_j F_4(r_j-r_k)}} \powe{S_2'},\\
\powe{S_2'} &= \powe{-\sum_{j<k} \klb*{A_jA_k \bra*{\phi(0)^2} + B_jB_k \bra*{\theta(0)^2} + A_jB_k \bra*{\phi(0)\theta(0)} + A_kB_j \bra*{\theta(0)\phi(0)}}}.
\end{align}
Combining $\powe{S_1}$ and $\powe{S_2'}$ with $\klb*{\phi(0),\theta(0)} = 0$ reduces to
\begin{align}
\powe{S_1+S_2'} &= \powe{-\frac{1}{2} \kla*{\sum_jA_j}^2 \bra*{\phi(0)^2} - \frac{1}{2} \kla*{\sum_jB_j}^2 \bra*{\theta(0)^2} - \kla*{\sum_jA_j}\kla*{\sum_jB_j} \bra*{\phi(0)\theta(0)}}.
\end{align}
Since both $\bra*{\phi(0)^2}$ and $\bra*{\theta(0)^2}$ diverge the full expectation value can only be non-zero if the so-called ``neutrality condition'' $\sum_j A_j = 0 = \sum_j B_j$ is met.
Under this constraint the result simplifies to
\begin{align}
\bra*{\prod_j \powe{\i A_j\phi(r_j) + \i B_j\theta(r_j)}} &= \powe{\frac{1}{2} \sum_{j<k} \klb*{A_jA_k F_1(r_j-r_k) + B_jB_kF_3(r_j-r_k) - A_jB_k F_2(r_j-r_k) - A_kB_j F_4(r_j-r_k)}}
\end{align}
and one only needs to compute the functions $F_1$ through $F_4$.
This result is independent of additional interaction terms in the Hamiltonian since we did not make any assumptions on the functions $F_j$.
For a quadratic Luttinger Hamiltonian it turns out that $F_4=F_2$ and $F_3 = \frac{1}{K^2} F_1$, further simplifying the formula.
\subsection{Path integrals via Fourier mode decomposition}
We can decompose the fields into Fourier series via\footnote{$r=(x,y) = (x,u\tau)$, $\textbf{q} = (k,\omega_n^\B / u)$ and $r\cdot \textbf{q} = kx - \omega_n^\B \tau$}
\begin{align}
\phi(r) &= \sum_\textbf{q} \powe{\i \textbf{q}\cdot r} \phi(\textbf{q})
\end{align}
and similarly for $\theta$.
Inserting these representations into the action in imaginary time yields
\begin{align}
-S &= \intx{0}{\beta}{\intr{\klb*{\frac{1}{\pi} \partial_x\theta \i\partial_\tau\phi - \mathcal{H}}}{x} }{\tau}\\
&= \sum_{\textbf{q},\textbf{q}'} \intx{0}{\beta}{\intr{\klb*{\frac{1}{\pi} \i k\theta(\textbf{q}) \omega_n'^\B \phi(\textbf{q}') - \frac{uK}{2\pi} \i k\theta(\textbf{q}) \i k'\theta(\textbf{q}') - \frac{u}{2\pi K} \i k\phi(\textbf{q}) \i k' \phi(\textbf{q}')} \powe{\i r\cdot (\textbf{q} + \textbf{q}')}}{x} }{\tau}\\
&= 2\pi\beta\sum_\textbf{q} \klb*{\frac{-\i k\omega_n^\B}{\pi} \theta(\textbf{q}) \phi^\ast(\textbf{q}) - \frac{uKk^2}{2\pi} \theta(\textbf{q}) \theta^\ast(\textbf{q}) - \frac{uk^2}{2\pi K} \phi(\textbf{q}) \phi^\ast(\textbf{q})}\\
&= -\beta \sum_\textbf{q} \Phi^\dagger(\textbf{q}) M(\textbf{q}) \Phi(\textbf{q})
\end{align}
where we defined the vector $\Phi(\textbf{q}) = \kla[\big]{\phi(\textbf{q}), \theta(\textbf{q})}^T$ and the matrix
\begin{align}
M &= \begin{pmatrix}
\frac{uk^2}{K} & \i k\omega_n^\B \\ \i k\omega_n^\B & uKk^2
\end{pmatrix}.
\end{align}
Including source terms we can then compute the partition function
\begin{align}
Z[J] &= \int \powe{-S[\Phi]} \powe{\sum_\textbf{q} J\cdot \Phi} \,\mathcal{D}\Phi = \mathcal{N} \powe{\frac{1}{4\beta} \sum_\textbf{q} J^\dagger(\textbf{q}) M^{-1}(\textbf{q}) J(\textbf{q})}.
\end{align}
Here the normalization factor $\mathcal{N}$ is independent of the source term $J$ and thus drops out when computing expectation values.
In fact, we can for example write
\begin{align}
\bra*{\phi(\textbf{q}_1) \phi(\textbf{q}_2)} &= \frac{1}{Z[J=0]} \frac{\delta^2}{\delta J_\phi(\textbf{q}_1) \delta J_\phi(\textbf{q}_2)} Z[J=0] = \frac{\delta^2}{\delta J_\phi(\textbf{q}_1) \delta J_\phi(\textbf{q}_2)} \log Z[J=0].
\end{align}
To this end we now drop the constant and explicitly evaluate
\begin{align}
\log Z[J] &= \frac{1}{4\beta} \sum_\textbf{q} J^\dagger(\textbf{q}) M^{-1}(\textbf{q}) J(\textbf{q}) \\
&= \frac{1}{4\beta} \sum_\textbf{q} \frac{uKk^2 J_\phi J_\phi^\ast + \frac{uk^2}{K} J_\theta J_\theta^\ast - \i k\omega_n^\B J_\phi J_\theta^\ast - \i k\omega_n^\B J_\phi^\ast J_\theta}{u^2k^4 + k^2\kla*{\omega_n^\B}^2} .
\end{align}
This leads to the expressions
\begin{align}
\bra*{\phi(\textbf{q}) \phi(\textbf{q}')} &= \frac{1}{4\beta} \frac{uKk^2}{u^2k^4 + k^2 \kla*{\omega_n^\B}^2} \frac{\delta }{\delta J_\phi(\textbf{q}')} 2J(-\textbf{q}) = \delta_{\textbf{q},-\textbf{q}'} \frac{uK}{2\beta} \frac{1}{u^2k^2 + \kla*{\omega_n^\B}^2},\\
\bra*{\theta(\textbf{q}) \theta(\textbf{q}')} &= \delta_{\textbf{q},-\textbf{q}'} \frac{u}{2\beta K} \frac{1}{u^2k^2 + \kla*{\omega_n^\B}^2},\\
\bra*{\phi(\textbf{q}) \theta(\textbf{q}')} &= \delta_{\textbf{q},-\textbf{q}'} \frac{-\i\omega_n^\B}{2\beta k} \frac{1}{u^2k^2 + \kla*{\omega_n^\B}^2} = \bra*{\theta(\textbf{q}) \phi(\textbf{q}')}.
\end{align}
To get the expectation values in real space and imaginary time we can write
\begin{align}
\bra*{\klb*{\phi(r) - \phi(0)}^2} &= \sum_{\textbf{q},\textbf{q}'} \kla*{\powe{\i \textbf{q}\cdot r} - 1} \kla*{\powe{\i \textbf{q}'\cdot r} - 1} \bra*{\phi(\textbf{q}) \phi(\textbf{q}')} \\
&=\sum_\textbf{q} \frac{uK}{2\beta} \frac{2 - 2\cos\kla*{\textbf{q}\cdot r}}{u^2k^2 + \kla*{\omega_n^\B}^2} = K \frac{u}{\beta} \sum_{n\in\mathbb{Z}} \intr{\frac{1 - \cos\kla*{kx - \omega_n^\B \tau}}{u^2k^2 + \kla*{\omega_n^\B}^2}}{k} \equiv KF_1(r),\\
\bra*{\klb*{\theta(r) - \theta(0)}^2} &= \frac{1}{K} F_1(r),\\
2\bra*{\phi(r)\theta(0)} &= 2\sum_{\textbf{q},\textbf{q}'} \powe{\i \textbf{q}\cdot r} \bra*{\phi(\textbf{q})\theta(\textbf{q}')} = \sum_{n\in\mathbb{Z}} \intr{ \powe{-\i\omega_n^\B\tau} \frac{\omega_n^\B}{\beta k} \frac{\sin(kx)}{u^2k^2 + \kla*{\omega_n^\B}^2}}{k} \\
&= 2\bra*{\theta(r) \phi(0)} \equiv F_2(r).
\end{align}
Note that there is an alternative route to evaluating these expectations values.
If we define the fields 
\begin{align}
\gamma_\ell &= K\theta - \ell \phi
\end{align}
then the action becomes diagonal since inserting $\theta = \frac{\gamma_\L + \Gamma_\R}{2K}$ and $\phi = \frac{\gamma_\L - \gamma_\R}{2}$ yield the new action
\begin{align}
-S &= -\beta \sum_\textbf{q} \begin{pmatrix}
\gamma_\R^\ast & \gamma_\L^\ast
\end{pmatrix} \frac{1}{2} \begin{pmatrix}
-1& \frac{1}{K} \\ 1 & \frac{1}{K}
\end{pmatrix} M(\textbf{q}) \frac{1}{2} \begin{pmatrix}
-1&1 \\ \frac{1}{K} & \frac{1}{K}
\end{pmatrix} \begin{pmatrix}
\gamma_\R \\ \gamma_\L
\end{pmatrix} \\
&= -\frac{\beta}{2K} \sum_\textbf{q} \begin{pmatrix}
\gamma_\R^\ast & \gamma_\L^\ast
\end{pmatrix} \begin{pmatrix}
uk^2 - \i\omega_n^\B k & 0 \\ 0 & uk^2 + \i\omega_n^\B k
\end{pmatrix} \begin{pmatrix}
\gamma_\R \\ \gamma_\L
\end{pmatrix}.
\end{align}
The relevant part of the partition function now becomes
\begin{align}
\log Z[J] &= \frac{K}{2\beta} \sum_\textbf{q} \sum_\ell \frac{\gamma_\ell\gamma_\ell^\ast}{uk^2 - \ell\i\omega_n^\B k}.
\end{align}
In this formulation there are only two non-zero expectation values in momentum-space
\begin{align}
\bra*{\gamma_\ell(\textbf{q}) \gamma_{\ell'}(\textbf{q}')} &= \delta_{\ell\ell'} \delta_{\textbf{q},-\textbf{q}'} \frac{K}{2\beta k} \frac{1}{uk - \ell\i\omega_n^\B}.
\end{align}
In real space we thus find the fundamental correlation functions
\begin{align}
\bra*{\gamma_\ell(r) \gamma_{\ell}(0)} &= \sum_{\textbf{q},\textbf{q}'} \powe{\i \textbf{q}\cdot r} \bra*{\gamma_\ell(\textbf{q}) \gamma_\ell(\textbf{q}')} = \frac{K}{2\beta} \sum_{n\in\mathbb{Z}} \intr{\powe{\i kx - \i\omega_n^\B\tau} \frac{1}{k} \frac{1}{uk - \ell\i\omega_n^\B}}{k}.
\end{align}
The integral diverges at $k=0$ so we instead compute
\begin{align}
\bra*{\gamma_\ell(r) \gamma_{\ell}(0) - \gamma_\ell(0)^2} &= \frac{K}{2\beta} \sum_{n\in\mathbb{Z}} \intr{\kla*{\powe{\i kx - \i\omega_n^\B\tau} - 1} \frac{1}{k} \frac{1}{uk - \ell\i\omega_n^\B}}{k}.
\end{align}
By symmetry $\bra*{\gamma_\R(x,\tau) \gamma_\R(0,0)} = \bra*{\gamma_\L(x,-\tau)\gamma_\L(0,0)}$ so we only need to compute one of the functions.
Using the summation formula\footnote{$\klb*{b} = b - \lfloor b \rfloor$ denotes the fractional part}
\begin{align}
\sum_{n\in\mathbb{Z}} \frac{\powe{2\pi \i n b}}{2\pi\i n + a} &= \frac{\powe{-a\klb*{b}}}{1 - \powe{-a}}
\end{align}
%a=-0.54;b=-0.001;n=np.arange(-4000,4001)
%np.sum(1/1j * np.exp(1j*b*(2*np.pi*n))/(2*np.pi*n-1j*a)), np.exp(a*(1 - b%1.0))/(np.exp(a)-1)
and separating the integral into regions with $k>0$ and $k<0$ reduces the Fourier transform to
\begin{align}
\bra*{\gamma_\L(r) \gamma_{\L}(0) - \gamma_\L(0)^2} &= \frac{K}{2} \intx{0}{\infty}{\frac{1}{k} \sum_{n\in\mathbb{Z}} \frac{\powe{-\alpha k}}{2\pi \i n + \beta uk} \klb*{\powe{\i kx - 2\pi\i n \frac{\tau}{\beta} } - 1} }{k} + (x,\tau\rightarrow -x,-\tau)\\
&= \frac{K}{2} \intx{0}{\infty}{\frac{1}{k} \frac{\powe{-\alpha k}}{1 - \powe{-\beta uk}} \klb*{\powe{\i kx + \beta uk \frac{\tau}{\beta}} - 1 }}{k} + (x,\tau\rightarrow -x,-\tau)\\
&= K \intx{0}{\infty}{\frac{1}{k} \frac{\powe{-\alpha k}}{1 - \powe{-\beta uk}} \klb*{\cosh\kla*{\i kx + uk\tau} - 1} }{k}\\
&= K\intx{0}{\infty}{\sum_{n\ge 0} \sum_{j\ge 1} \powe{-(\beta un + \alpha) k} \frac{1}{(2j)!} (\i x+u\tau)^{2j} k^{2j-1} }{k}\\
&= K \sum_{n\ge 0} \sum_{j\ge 1} \frac{1}{(2j)!} \kla*{\frac{\i x+u\tau}{\beta un + \alpha}}^{2j} (2j-1)!\\
&= -K \sum_{n\ge 0} \log\kla*{1 - \frac{(\i x+u\tau)^2}{(\beta un + \alpha)^2}}\\
&\simeq K\log\kla*{1 - \frac{(\i x + u\tau)^2}{\alpha^2}} + K \log\kla*{\prod_{n\ge 1} \frac{(\beta un)^2}{(\beta un)^2 + (x - \i u\tau)^2}} \\
&= K\log\kla*{1 - \frac{(\i x + u\tau)^2}{\alpha^2}} + K \log\kla*{\Gamma\kla*{1 + \frac{\i x + u\tau}{\beta u}} \Gamma\kla*{1 - \frac{\i x + u\tau}{\beta u}}}\\
&= K\log\kla*{1 - \frac{(\i x + u\tau)^2}{\alpha^2}} + K \log \kla*{\frac{\frac{\pi}{\beta u} \klb*{u\tau + \i x}}{\sin\kla*{\frac{\pi}{\beta u} \klb*{u\tau + \i x}}}}.
\end{align}
%By symmetry $\bra*{\gamma_\R(x,\tau) \gamma_\R(0,0)} = \bra*{\gamma_\L(x,-\tau)\gamma_\L(0,0)}$ so we only need to compute one of the functions.
%Also note that we can replace $x$ by $|x|$ and $\tau$ by $\tau\sgn(x)$ by substituting $k\rightarrow -k$ and $\omega_n^\B\rightarrow -\omega_n^\B$.
%We will regularize the integral by inserting a small parameter $\delta$ in the denominator.
%The integrand still diverges at $k=0$ though, and we therefore need to evaluate the integral around a small contour parametrized by $k(t) = \epsilon\powe{\i t}$.
%This gives
%\begin{align}
%\lim_{\epsilon\rightarrow 0} \frac{K}{2\beta} \sum_{n\in\mathbb{Z}} \powe{-\i\omega_n^\B\tau \sgn(x)} \intx{\pi}{0}{\powe{\i |x| \epsilon\powe{\i t}} \frac{1}{\epsilon\powe{\i t}} \frac{1}{u\epsilon\powe{\i t} + \i\omega_n^\B + \delta} \epsilon\i\powe{\i t}}{t} &= \frac{K}{2\beta} \sum_{n\in\mathbb{Z}} \powe{-\i\omega_n^\B \tau\sgn(x)} \frac{-\pi\i}{\i\omega_n^\B + \delta}.
%\end{align}
%Closing the contour in the upper half of the complex plane then yields
%\begin{align}
%\bra*{\gamma_\L(r) \gamma_\L(0)} &= \frac{K}{2\beta} \sum_{n\in\mathbb{Z}} \powe{-\i\omega_n^\B\tau\sgn(x)} \klb*{ \frac{\pi\i}{\i\omega_n^\B + \delta} + 2\pi\i \powe{\i k|x|} \frac{1}{k} \Big\vert_{uk=-\i\omega_n^\B - \delta} \Theta(-n)}•
%\end{align}
%Using the summation formula\footnote{$\klb*{b} = b - \lfloor b \rfloor$ denotes the fractional part}
%\begin{align}
%\sum_{n\in\mathbb{Z}} \frac{\powe{2\pi \i n b}}{2\pi\i n + a} &= \frac{\powe{-a\klb*{b}}}{1 - \powe{-a}}
%\end{align}
\section{Interacting Luttinger Hamiltonian}
Consider the quadratic Luttinger Hamiltonian with an additional potential of the form
\begin{align}
H &= \frac{u}{2\pi} \intr{\klb*{K(\partial_x\theta)^2 + \frac{1}{K} (\partial_x\phi)^2}}{x} + g\intr{V(\phi)}{x}.
\end{align}
Using $\partial_t\phi = uK\partial_x\theta = uK\pi\Pi$ and $t=-\i\tau = \frac{y}{\i u}$ with $b=\beta u$ yields the action
\begin{align}
S &= \intx{0}{\beta}{\intr{\klb*{\Pi \partial_t\phi - \frac{uK}{2\pi} \frac{(\partial_t\phi)^2}{u^2K^2} - \frac{u}{2\pi K} (\partial_x\phi)^2 - gV}}{x}}{\tau} \\
&= -\frac{1}{2\pi K} \intx{0}{b}{\intr{\klb*{(\partial_y\phi)^2 + (\partial_x\phi)^2 + \tilde{g} V}}{x} }{y}
\end{align}
with the rescaled coupling $\tilde{g} = g\frac{2\pi K}{u}$.
As a semi-classical approximation we will expand the potential around the classical solution by writing $\phi=\phi_\text{cl} + \delta\phi$ giving
\begin{align}
S &= \bsplit{S_\text{cl} - \frac{1}{2\pi K}\intx{0}{b}{\intr{\klb[\Big]{&2\partial_y\phi_\text{cl} \partial_\tau\delta\phi + (\partial_y\delta\phi)^2 + 2\partial_x\phi_\text{cl} \partial_x\delta\phi + (\partial_x\delta\phi)^2 \\
&+ \tilde{g} V'(\phi_\text{cl}) \delta\phi + \frac{\tilde{g}}{2} V''(\phi_\text{cl}) \delta\phi^2}}{x} }{y}}\\
&= S_\text{cl} - \frac{1}{2\pi K} \intx{0}{b}{\intr{\klb*{(\partial_y\delta\phi)^2 + (\partial_x\delta\phi)^2 + \frac{\tilde{g}}{2} V''(\phi_\text{cl})\delta\phi^2}}{x} }{y}.
\end{align}
Here the first order term vanished due to the classical equation of motion
\begin{align}
0 &= \partial_\mu \frac{\partial\mathcal{L}}{\partial \partial_\mu\phi} - \frac{\partial\mathcal{L}}{\partial\phi} = \partial_y \frac{-u}{\pi K} \partial_y\phi + \partial_x \frac{-u}{\pi K} \partial_x\phi - \frac{-u\tilde{g}}{2\pi K} V'(\phi).
\end{align}
Specifying $\phi=\phi_\text{cl}$ we find
\begin{align}
\kla*{\partial_y^2 + \partial_x^2}\phi_\text{cl} &= \frac{\tilde{g}}{2} V'(\phi_\text{cl}).
\end{align}
For now let us assume that we already have a solution $\phi_\text{cl}$ to this equation.
Then the path integral can be written as
\begin{align}
Z &= Z_\text{cl} \int \exp\kla*{\frac{1}{2\pi K} \intx{0}{b}{\intr{\klb*{(\partial_y\delta\phi)^2 + (\partial_x\delta\phi)^2 + \frac{\tilde{g}}{2} V''(\phi_\text{cl}) \delta\phi^2}}{x} }{y}}\,\mathcal{D}\delta\phi
\end{align}
where the classical part is $Z_\text{cl} = \powe{-S_\text{cl}}$.
Since this part will always cancel when calculating expectation values, we will drop it from now on.
For example
\begin{align}
\bra*{\phi(r)\phi(0)} &= \bra*{\phi_\text{cl}(r)\phi_\text{cl}(0) + \phi_\text{cl}(r)\delta\phi(0) + \phi_\text{cl}(0)\delta\phi(r) + \delta\phi(r)\delta\phi(0)} = \\
&= \phi_\text{cl}(r)\phi_\text{cl}(0) + \bra*{\delta\phi(r)\delta\phi(0)}
\end{align}
since the action is quadratic in $\delta\phi$ and therefore $\bra*{\delta\phi} = 0$ by symmetry.
Note that in the non-interacting case the simplest classical solution is $\phi_\text{cl}=0$, and thus we recover the path integral of the free model by replacing $\delta\phi$ with $\phi$.\\
The deviation must vanish at the boundaries which allows us to integrate by parts without picking up additional terms.
This results in
\begin{align}
Z[J] &= \int \exp\kla*{-\frac{1}{2\pi K} \intx{0}{b}{\intr{\delta\phi \klb*{\partial_y^2 + \partial_x^2 - \frac{\tilde{g}}{2} V''(\phi_\text{cl})} \delta\phi}{x} }{y}} \powe{\int J(r)\delta\phi}\,\mathcal{D}\delta\phi\\
&= \mathcal{N} \powe{\frac{\pi K}{2} \int J(r)D(r-r')J(r')}
\end{align}
where $\mathcal{N}$ is some normalization independent of the source term $J$ and the kernel satisfies
\begin{align}
\kla*{\partial_x^2 + \partial_y^2 - \frac{\tilde{g}}{2} V''(\phi_\text{cl})} D(r) &= \delta(r).
\end{align}
Given the kernel we can immediately write down the expectation value
\begin{align}
\bra*{\klb*{\phi(r) - \phi(0)}^2} &= -2\bra*{\phi(r)\phi(0) - \phi(0)^2} = -2\frac{\delta^2}{\delta J(r)\delta J(0)} \log Z[J] + 2\frac{\delta^2}{\delta J(0)^2}\log Z[J] \\
&= -2\pi K\klb*{D(r) - D(0)}.
\end{align}
In contrast to the free case the value of $D(0)$ is neither irrelevant nor arbitrary but rather dictated by the interaction potential.
\subsection{Classical solution for the sine-Gordon potential}
Let us focus on the sine-Gordon potential $V(\phi) = \cos(m\phi)$.
Then we want to find a solution to the equation
\begin{align}
\nabla^2\phi_\text{cl} &= -\frac{\tilde{g} m}{2} \sin\kla*{m\phi_\text{cl}}
\end{align}
that satisfies a periodic boundary condition $\phi_\text{cl}(y+b) = \phi_\text{cl}(y)$.
Defining $s(x,y) = m\phi_\text{cl}(x,y)$ yields the typical form of the sine-Gordon equation in imaginary time
\begin{align}
\nabla^2s &= -\frac{\tilde{g} m^2}{2} \sin (s) \equiv M^2 \sin(s).
\end{align}
This equation can be solved by $s(x,y) = 4\arctan\kla*{\powe{c_0 + M\gamma \kla*{x \pm v\i y}}}$ where $\gamma = \frac{1}{\sqrt{1 - v^2}}$, $M = \i m\sqrt{\frac{\tilde{g}}{2}}$ with arbitrary parameters $c_0$ and $v$.
%c1,c2,c3,c4,x,y,m,b = sy.symbols("c_1 c_2 c_3 c_4 x y m b", real=True)
%phi = c1 *sy.atan(c2 * sy.exp(c3 * x + c4 * y))
%pot = sy.cos(m*phi)
%pot_prime = -m * sy.sin(m*phi)
%phi_xx = sy.diff(phi, x, 2)
%phi_yy = sy.diff(phi, y, 2)
%print((phi_xx+phi_yy).simplify())
%%%
%###real time
%s = sy.simplify(4*sy.atan(c2 * sy.exp(1/sy.sqrt(1-c3**2) * m * (x - c3*y))))
%rhs = sy.trigsimp(m**2*sy.sin(s))
%lhs = sy.simplify(sy.diff(s,x,2)-sy.diff(s,y,2))
%(rhs-lhs).simplify()
%###imag time
%s = sy.simplify(4*sy.atan(c2 * sy.exp(1/sy.sqrt(1-c3**2) * m * (x - sy.I* c3*y))))
%rhs = sy.trigsimp(m**2*sy.sin(s))
%lhs = sy.simplify(sy.diff(s,x,2)+sy.diff(s,y,2))
%(rhs-lhs).simplify()
In order to satisfy the periodic boundary condition we require
\begin{align}
1 &= \powe{\pm M\gamma v\i b} = \powe{\mp m \sqrt{\frac{\tilde{g}}{2}} \frac{v}{\sqrt{1 - v^2}} b}.
\end{align}
This leads to the condition $\frac{2\pi\i}{mb} \sqrt{\frac{2}{\tilde{g}}} = \frac{v}{\sqrt{1 - v^2}}$ thus\footnote{$v= c\sqrt{1 - v^2}$ has the solutions $v = \pm \frac{c}{\sqrt{1 + c^2}}$}
\begin{align}
v &= \pm \klb*{1 + \kla*{\frac{mb}{2\pi\i }}^2 \frac{2}{\tilde{g}}}^{-\frac{1}{2}}.
\end{align}
Choosing $c_0=0$ we find
\begin{align}
\phi_\text{cl}(x,y) &= \frac{4}{m} \arctan\kla*{\powe{\i \frac{2\pi}{b} y} \powe{\i m \sqrt{\frac{\tilde{g}}{2}} \gamma x}}.
\end{align}
However, note that the trivial solution $\phi_\text{cl} = 0$ has a classical action that diverges toward negative infinity.
\subsection{Semi-classical approximation for the trivial classical solution}
For now we will hence restrict our analysis to the simplest case of $\phi_\text{cl}\equiv 0$ as it should bring the largest contribution to the path integral.
The kernel equation then becomes
\begin{align}
\kla*{\nabla^2 + \Lambda^2}D(r) &= \delta(r)
\end{align}
%%% HELMHOLTZ equation
%subs = {x : 0.32, y : 0.22, b : 7}
%f = -sy.I/4 * sy.hankel1(0,b * sy.sqrt(x**2+y**2))
%deriv = (f.diff(x,2) + f.diff(y,2) + b**2 * f).simplify()
%print(deriv.evalf(subs=subs))
%f = -sy.I/4 * sy.hankel1(0,b * sy.sqrt(sy.sinh(x)**2+sy.sin(y)**2))
%deriv = (f.diff(x,2) + f.diff(y,2) + b**2 * f).simplify()
%print(deriv.evalf(subs=subs))
%%% doesn't work though, need a way to adjust this solution to PBC...
where we introduced the constant $\Lambda^2=\frac{\tilde{g}}{2} m^2$.
Completely analogous to the calculation in the free case we Fourier transform to get the algebraic equation
\begin{align}
1 &= \kla*{\Lambda^2 - w_n^2 - k^2}D(k,w_n)
\end{align}
with the discrete frequencies $w_n = \frac{2\pi n}{b}$ for $n\in\mathbb{Z}$ and evaluate the inverse transform
\begin{align}
D(r) &= \frac{1}{2\pi b} \sum_{n\in\mathbb{Z}} \powe{-\i w_ny} \intr{\powe{\i kx} \frac{1}{\Lambda^2 - w_n^2 - k^2}}{k}\\
&= \frac{1}{2\pi b} \sum_{n\in\mathbb{Z}} \powe{-\i w_n|y|} \intr{\powe{\i k|x|} \frac{-1}{\kla*{k - \i\sqrt{w_n^2 - \Lambda^2}} \kla*{k + \i\sqrt{w_n^2 - \Lambda^2}}}}{k}.
\end{align}
Closing the contour in the upper half of the complex plane results in the residues $k_n = \i\sqrt{w_n^2 - \Lambda^2}$ lying inside the contour for $|n| > \frac{\Lambda b}{2\pi}$.
Note that the other poles are pushed to the real axis and thus need to be avoided.
Let $|n| < \frac{\Lambda b}{2\pi}$.
Then integrating in a small semi-circle with the parametrization $k_n + \epsilon \powe{\i t}$ yields
\begin{align}
\lim_{\epsilon\rightarrow 0}\intx{\pi}{0}{\powe{\i |x| k_n + \i |x| \epsilon \powe{\i t}} \frac{-1}{\epsilon \powe{\i t} \kla*{2k_n + \epsilon\powe{\i t}}} \epsilon\i \powe{\i t}}{t} &= \frac{\pi\i}{2k_n} \powe{\i |x|k_n}.
\end{align}
Therefore we find that the integral decomposes into two contributions
\begin{align}
D(r) &= \frac{1}{2\pi b} \sum_{|n|< \frac{\Lambda b}{2\pi}} \powe{-\i w_n |y|} \frac{-\pi\i}{2k_n} \powe{\i |x| k_n} + \frac{1}{2\pi b} \sum_{|n|> \frac{\Lambda b}{2\pi}} \powe{-\i w_n |y|} (-2\pi\i) \frac{1}{2k_n} \powe{\i |x| k_n}\\
&= \frac{1}{4\i b} \sum_{|n|<\frac{\Lambda b}{2\pi}} \frac{1}{\sqrt{\Lambda^2 - w_n^2}} \powe{\i |x|\sqrt{\Lambda^2 - w_n^2} - \i w_n |y|} - \frac{1}{2 b} \sum_{|n|> \frac{\Lambda b}{2\pi}} \frac{1}{\sqrt{w_n^2 - \Lambda^2}} \powe{-|x|\sqrt{w_n^2 - \Lambda^2} - \i w_n |y|}\\
&= \frac{1}{4\i b\Lambda} \powe{\i |x|\Lambda} + \frac{1}{2\i b} \sum_{0<n<\frac{\Lambda b}{2\pi}} \frac{\cos(w_n y)}{\sqrt{\Lambda^2 - w_n^2}} \powe{\i|x|\sqrt{\Lambda^2 - w_n^2}} - \frac{1}{b} \sum_{n> \frac{\Lambda b}{2\pi}} \frac{\cos(w_n y)}{\sqrt{w_n^2 - \Lambda^2}} \powe{-|x|\sqrt{w_n^2 - \Lambda^2}}.
\end{align}
%b=1.4
%lam=7
%n=2
%omega=2*np.pi*n/b
%x=0.59
%y=0.43
%##integrate cos(0.59 x) / (-31.568199 - x^2) from x=-inf to inf
%numeric = integrate.quad(lambda k: np.cos(np.abs(x)*k) / (lam**2 - omega**2 - k**2), -np.inf, np.inf, limit=10000)
%if np.abs(n) > lam * b / (2*np.pi):
%    analytic = -np.pi * np.exp(-np.abs(x) * np.sqrt(omega**2 - lam**2)) / np.sqrt(omega**2 - lam**2)
%elif n == 0:
%    analytic = np.pi * np.exp(1j*np.abs(x) * lam) / (2*1j*lam)
%else:
%    analytic = np.pi/(2*1j)* np.exp(1j*np.abs(x) * np.sqrt(lam**2 - omega**2)) / np.sqrt(lam**2 - omega**2)
\subsection{Path integrals via Fourier mode decomposition}
We use the semi-classical approximation $\cos(\phi) \simeq -\frac{1}{2} \delta\phi^2$ for the sine-Gordon potential.
As argued above, we may write $\phi$ instead of $\delta\phi$ since the classical part vanishes in any expression anyway.
We can thus once again decompose the fields into Fourier series via
\begin{align}
\phi(r) &= \sum_\textbf{q} \powe{\i \textbf{q}\cdot r} \phi(\textbf{q})
\end{align}
and similarly for $\theta$.
Inserting these representations into the action in imaginary time yields
\begin{align}
-S &= \intx{0}{\beta}{\intr{\klb*{\frac{1}{\pi} \partial_x\theta \i\partial_\tau\phi - \mathcal{H}}}{x} }{\tau}\\
&\simeq \sum_{\textbf{q},\textbf{q}'} \intx{0}{\beta}{\intr{\klb*{\frac{1}{\pi} \i k\theta(\textbf{q}) \omega_n'^\B \phi(\textbf{q}') - \frac{uK}{2\pi} \i k\theta(\textbf{q}) \i k'\theta(\textbf{q}') - \frac{u}{2\pi K} \i k\phi(\textbf{q}) \i k' \phi(\textbf{q}') + \frac{g}{2} \phi(\textbf{q}) \phi(\textbf{q}')} \powe{\i r\cdot (\textbf{q} + \textbf{q}')}}{x} }{\tau}\\
&= 2\pi\beta\sum_\textbf{q} \klb*{\frac{-\i k\omega_n^\B}{\pi} \theta(\textbf{q}) \phi^\ast(\textbf{q}) - \frac{uKk^2}{2\pi} \theta(\textbf{q}) \theta^\ast(\textbf{q}) - \frac{uk^2}{2\pi K} \phi(\textbf{q}) \phi^\ast(\textbf{q}) + \frac{g}{2} \phi(\textbf{q}) \phi^\ast(\textbf{q})}\\
&= -\beta \sum_\textbf{q} \Phi^\dagger(\textbf{q}) M(\textbf{q}) \Phi(\textbf{q})
\end{align}
where we defined the vector $\Phi(\textbf{q}) = \kla[\big]{\phi(\textbf{q}), \theta(\textbf{q})}^T$ and the matrix
\begin{align}
M &= \begin{pmatrix}
\frac{uk^2}{K} - \pi g & \i k\omega_n^\B \\ \i k\omega_n^\B & uKk^2
\end{pmatrix}.
\end{align}
Including source terms we can then compute the partition function
\begin{align}
Z[J] &= \int \powe{-S[\Phi]} \powe{\sum_\textbf{q} J\cdot \Phi} \,\mathcal{D}\Phi = \mathcal{N} \powe{\frac{1}{4\beta} \sum_\textbf{q} J^\dagger(\textbf{q}) M^{-1}(\textbf{q}) J(\textbf{q})}.
\end{align}
Here the normalization factor $\mathcal{N}$ is independent of the source term $J$ and thus drops out when computing expectation values.
Up to the aforementioned constant the logarithm of the partition function thus becomes
\begin{align}
\log Z[J] &= \frac{1}{4\beta} \sum_\textbf{q} J^\dagger(\textbf{q}) M^{-1}(\textbf{q}) J(\textbf{q}) \\
&= \frac{1}{4\beta} \sum_\textbf{q} \frac{uKk^2 J_\phi J_\phi^\ast + \kla*{\frac{uk^2}{K} - \pi g} J_\theta J_\theta^\ast - \i k\omega_n^\B J_\phi J_\theta^\ast - \i k\omega_n^\B J_\phi^\ast J_\theta}{uKk^2 \kla*{\frac{uk^2}{K} - \pi g} + k^2\kla*{\omega_n^\B}^2} .
\end{align}
This leads to the expressions
\begin{align}
\bra*{\phi(\textbf{q}) \phi(\textbf{q}')} &= \frac{1}{4\beta} \frac{uKk^2}{uKk^2 \kla*{\frac{uk^2}{K} - \pi g} + k^2\kla*{\omega_n^\B}^2} \frac{\delta }{\delta J_\phi(\textbf{q}')} 2J(-\textbf{q}) = \delta_{\textbf{q},-\textbf{q}'} \frac{uK}{2\beta} \frac{1}{u^2k^2 - \pi ugK + \kla*{\omega_n^\B}^2},\\
\bra*{\theta(\textbf{q}) \theta(\textbf{q}')} &= \delta_{\textbf{q},-\textbf{q}'} \frac{1}{2\beta} \frac{\frac{u}{K} - \pi \frac{g}{k^2}}{u^2k^2 - \pi ugK + \kla*{\omega_n^\B}^2},\\
\bra*{\phi(\textbf{q}) \theta(\textbf{q}')} &= \delta_{\textbf{q},-\textbf{q}'} \frac{-\i\omega_n^\B}{2\beta k} \frac{1}{u^2k^2 - \pi ugK + \kla*{\omega_n^\B}^2} = \bra*{\theta(\textbf{q}) \phi(\textbf{q}')}.
\end{align}
To get the expectation values in real space and imaginary time we can write
\begin{align}
\bra*{\klb*{\phi(r) - \phi(0)}^2} &= \sum_{\textbf{q},\textbf{q}'} \kla*{\powe{\i \textbf{q}\cdot r} - 1} \kla*{\powe{\i \textbf{q}'\cdot r} - 1} \bra*{\phi(\textbf{q}) \phi(\textbf{q}')} \\
&= K \frac{u}{\beta} \sum_{n\in\mathbb{Z}} \intr{\frac{1 - \cos\kla*{kx - \omega_n^\B \tau}}{u^2k^2 + \kla*{\omega_n^\B}^2 - \pi ugK}}{k} \equiv KF_1(r),\\
\bra*{\klb*{\theta(r) - \theta(0)}^2} &= \frac{1}{K}F_1(r) - \frac{\pi g}{\beta} \sum_{n\in\mathbb{Z}} \intr{ \frac{1}{k^2} \frac{1 - \cos\kla*{kx - \omega_n^\B \tau}}{u^2k^2 + \kla*{\omega_n^\B}^2 - \pi ugK}}{k} \equiv \frac{1}{K} F_1(r) - gu\beta^2 F_3(r),\\
2\bra*{\phi(r)\theta(0)} &= 2\sum_{\textbf{q},\textbf{q}'} \powe{\i \textbf{q}\cdot r} \bra*{\phi(\textbf{q})\theta(\textbf{q}')} = \sum_{n\in\mathbb{Z}} \intr{\powe{-\i\omega_n^\B\tau} \frac{\omega_n^\B}{\beta k} \frac{\sin(kx)}{u^2k^2 + \kla*{\omega_n^\B}^2 - \pi ugK}}{k} \\
&= 2\bra*{\theta(r) \phi(0)} \equiv F_2(r).
\end{align}
The special functions we have to compute are
\begin{align}
F_1(r) &= \frac{u}{\beta} \sum_{n\in\mathbb{Z}} \intr{\frac{1 - \cos\kla*{kx - \omega_n^\B\tau}}{u^2k^2 + \kla*{\omega_n^\B}^2 - \pi ugK}}{k},\\
F_2(r) &= \frac{1}{\beta} \sum_{n\in\mathbb{Z}} \intr{\powe{-\i\omega_n^\B\tau} \frac{\omega_n^\B}{k} \frac{\sin(kx)}{u^2k^2 + \kla*{\omega_n^\B}^2 - \pi ugK}}{k},\\
F_3(r) &= \frac{\pi}{u\beta^3} \sum_{n\in\mathbb{Z}} \intr{ \frac{1}{k^2} \frac{1 - \cos\kla*{kx - \omega_n^\B \tau}}{u^2k^2 + \kla*{\omega_n^\B}^2 - \pi ugK}}{k}.
\end{align}
There is one more thing we have to consider before diving into the calculations.
In the general formula for expectation values of products of vertex operators we assumed $\bra*{\phi(0)^2}$ and $\bra*{\theta(0)^2}$ to diverge towards infinity.
However, due to the interaction term we instead find
\begin{align}
\bra*{\phi(0)^2} &= \sum_{\textbf{q},\textbf{q}'} \bra*{\phi(\textbf{q})\phi(\textbf{q}')} = \frac{uK}{2\beta} \sum_{n\in\mathbb{Z}} \intr{\frac{1}{u^2k^2 + \kla*{\omega_n^\B}^2 - \pi ugK}}{k} = ???
\end{align}



%\begin{thebibliography}{9}
%\bibitem{example} description, \url{https://www.exmapleurl}, day.\,month~2021.
%\end{thebibliography}

\end{document}